\chapter{Квазиньютоновские методы и мультисекущие обновления}\label{ch:jacobian}

Большинство задач численной оптимизации и~решения нелинейных систем
сводится к~поиску точки, в~которой обращается в~нуль градиент~$\nabla f$
или нелинейное отображение~$F$. Простейший подход~--- методы первого
порядка, прежде всего градиентный спуск
$x_{k+1}=x_k-\alpha\,\nabla f(x_k)$~\cite{polyak1983,nesterov2010}: он использует только градиент
и~дёшев на~итерации, но~не~учитывает кривизну и~потому сходится
линейно, с~фактором, определяемым обусловленностью задачи и~растущим
вместе с~ней. Метод Ньютона учитывает кривизну через
якобиан (для~минимизации~--- гессиан) и~сходится локально квадратично,
однако требует вычисления этой производной и~решения линейной системы
с~ней на~каждом шаге, что~для крупных или сложно дифференцируемых
задач неприемлемо дорого. Квазиньютоновские методы занимают
промежуточное положение: они заменяют производную дешёвой
аппроксимацией~$B_k$, обновляемой по~ходу итераций из~уже вычисленных
значений~$F$, и~сохраняют сверхлинейную сходимость без~явного
дифференцирования. Настоящая глава вводит аппарат этих методов
и~помещает работу в~контекст известных результатов: от~вывода
квазиньютоновской итерации и~классических обновлений до~общей
мультисекущей формулы, на~которой строится анализ последующих глав.

%% ============================================================
\section{Задачи оптимизации и нелинейные системы}\label{sec:ch1_setup}
%% ============================================================

Рассмотрим задачу безусловной минимизации гладкой функции
\begin{equation}\label{eq:min}
  \min_{x\in\RR^n} f(x), \qquad f\colon\RR^n\to\RR,\;\; f\in C^2.
\end{equation}
Частным примером этой задачи является оценка параметров моделей и обучение нейросетей; систематическое изложение численных методов оптимизации дано в~\cite{izmailov2008}. В~точке минимума выполнено условие стационарности
$\nabla f(x)=0$~--- система $n$ уравнений с~$n$ неизвестными, частный
случай более общей задачи поиска корня нелинейного отображения
\begin{equation}\label{eq:system}
  F(x) = 0, \qquad F\colon \RR^n \to \RR^n.
\end{equation}
При $F=\nabla f$ задача~\eqref{eq:system} есть условие стационарности
для~\eqref{eq:min}; в~общем случае $F$ возникает самостоятельно~---
при дискретизации краевых задач, в~равновесных и~экономических моделях,
а~также как оператор оптимальности седловых (минимаксных) задач
$\min_x\max_y \Phi(x,y)$, где $F(z)=(\nabla_x\Phi,\,-\nabla_y\Phi)$,
$z=(x,y)$, и~вариационных неравенств~(ВН)~\cite{viqa2024}.

Структурное различие между этими двумя источниками~--- симметрия
производной. Для~\eqref{eq:min} матрицей системы $\nabla f(x)=0$ служит
гессиан $\nabla^2 f(x)$, симметричный по~построению; для общей
системы~\eqref{eq:system} якобиан $J(x)=F'(x)$ в~общем случае
несимметричен (таблица~\ref{tab:ch1_classes}).

\begin{table}[ht]
\centering
\caption{Два класса задач, структура производной и~отвечающие им
квазиньютоновские обновления.}\label{tab:ch1_classes}
\renewcommand{\arraystretch}{1.3}
\begin{tabular}{@{}p{4.4cm} p{3.9cm} p{4.4cm}@{}}
\toprule
Класс задачи & Производная & Обновления \\
\midrule
Оптимизация $\min f$,\newline $\nabla f(x)=0$
  & гессиан $\nabla^2 f$~--- симметричен
  & симметричные: SR1, PSB, DFP, BFGS (гл.~\ref{ch:symmetric}) \\
Система $F(x)=0$\newline (вкл.\ минимакс, ВН)
  & якобиан $F'$~--- несимметричен
  & несимметричные: Бройден (гл.~\ref{ch:sp-broyden-theory}) \\
\bottomrule
\end{tabular}
\end{table}


%% ============================================================
\section{От метода Ньютона к квазиньютоновским обновлениям}\label{sec:ch1_newton}
%% ============================================================

Метод Ньютона для системы~\eqref{eq:system} строит шаг из~линеаризации
$F(x_k+d)\approx F(x_k)+J(x_k)\,d$, приравнивая правую часть нулю:
$J(x_k)\,d_k=-F(x_k)$. Квазиньютоновская итерация заменяет в~нём
якобиан~$J(x_k)$ аппроксимацией~$B_k$, не~требующей явного
дифференцирования:
\begin{equation}\label{eq:step}
  x_{k+1} = x_k + d_k, \qquad B_k\,d_k = -F(x_k),
\end{equation}
где $B_k \in \RR^{n \times n}$~--- текущая аппроксимация якобиана
$J(x_k) = F'(x_k)$ (для~\eqref{eq:min}~--- гессиана $\nabla^2 f(x_k)$). Вырожденный
выбор $B_k=\alpha^{-1}I$ возвращает градиентный спуск без~информации
о~кривизне, тогда как $B_k=J(x_k)$~--- метод Ньютона.

Условие, которому подчиняется обновление $B_k\mapsto B_{k+1}$, проще
всего получить из~квадратичной модели минимизации~\cite[\S6.1]{bfgs}.
Приблизим $f$ вблизи $x_{k+1}$ моделью
\[
  m_{k+1}(d) = f(x_{k+1}) + \nabla f(x_{k+1})^\top d
               + \tfrac12\, d^\top B_{k+1}\, d
\]
и~потребуем, чтобы её градиент совпадал с~истинным не~только в~$d = 0$
(что~выполнено автоматически), но~и~в~предыдущей точке:
$\nabla m_{k+1}(-s_k)=\nabla f(x_k)$, где
\[
  s_k = x_{k+1} - x_k, \qquad y_k = \nabla f(x_{k+1}) - \nabla f(x_k).
\]
Подстановка $\nabla m_{k+1}(d)=\nabla f(x_{k+1})+B_{k+1}d$ при $d=-s_k$
даёт \emph{секущее условие}
\begin{equation}\label{eq:secant}
  B_{k+1}\, s_k = y_k.
\end{equation}
Если от~$B_{k+1}$ дополнительно требуется симметрия и~положительная
определённость (оптимизационный случай $F=\nabla f$, где $B_{k+1}$
аппроксимирует гессиан), то~секущее условие~\eqref{eq:secant} разрешимо
лишь при выполнении \emph{условия кривизны}
\begin{equation}\label{eq:curvature}
  s_k^\top y_k > 0,
\end{equation}
получаемого домножением~\eqref{eq:secant} слева на~$s_k^\top$: при
$B_{k+1}\succ 0$ имеем $s_k^\top y_k = s_k^\top B_{k+1} s_k > 0$. Для
сильно выпуклой~$f$ неравенство выполнено на~любой паре точек,
а~в~общем случае его гарантирует линейный поиск с~условием
Вольфа~\cite{wolfe1969} (\S\ref{sec:ch1_global}).
Для общей системы~\eqref{eq:system} та~же пара определяется через
$y_k = F(x_{k+1}) - F(x_k)$, и~условие~\eqref{eq:secant} выражает
требование, чтобы $B_{k+1}$ воспроизводило конечную разность $F$ вдоль
последнего шага~--- многомерный аналог секущего метода.

По~формуле
Тейлора с~остатком в~интегральной форме вдоль отрезка $[x_k,x_{k+1}]$
\begin{equation}\label{eq:avg_operator}
  y_k = \Bigl(\int_0^1 \nabla^2 f(x_k + t\,s_k)\,dt\Bigr) s_k
      = \bar G_k\, s_k,
\end{equation}
а в~общем случае $y_k = \bar J_k\, s_k,\quad \bar J_k = \int_0^1 J(x_k + t\,s_k)\,dt$, где $\bar G_k$~--- усреднённый вдоль шага гессиан, $\bar J_k$~---
усреднённый якобиан~\cite[\S6.1]{bfgs}. В~этих обозначениях секущее
условие~\eqref{eq:secant} принимает вид $B_{k+1} s_k = \bar J_k s_k$:
обновление обязано воспроизводить усреднённую производную \emph{ровно
вдоль сделанного шага}, оставаясь свободным на~ортогональном дополнении
к~$s_k$.

В главах~\ref{ch:sp-broyden-theory}--\ref{ch:viqa} отклонение усреднённого оператора~$\bar J_k$
от~якобиана в~решении $\Jstar = J(x^{*})$ управляет оценкой ограниченного
ухудшения (\emph{bounded deterioration}): в~доказательстве
Теоремы~\ref{thm:bd} (\S\ref{sec:sp-broyden-proof}) столбцы секущей
невязки
\[
  y_j - \Jstar s_j
  = \Bigl(\int_0^1 \bigl[J(x_j + t\,s_j) - \Jstar\bigr]\,dt\Bigr) s_j
\]
оцениваются через липшицеву норму подынтегрального
выражения~\eqref{eq:integral_bound}, и~эта~же конструкция лежит в~основе
анализа секуще-сохраняющих обновлений
главы~\ref{ch:viqa}.

Поскольку
\eqref{eq:secant}~--- это $n$ уравнений на~$n^2$ элементов $B_{k+1}$,
обновление недоопределено; разные способы доопределения порождают
семейство классических обновлений (\S\ref{sec:ch1_classical}) и~его
мультисекущие расширения (\S\ref{sec:ch1_multisecant}).
\section{Классические квазиньютоновские обновления}\label{sec:ch1_classical}
%% ============================================================

Классические обновления различаются тем, как доопределяется
недоопределённая секущая система~\eqref{eq:secant} (симметрия, ранг
коррекции, выбор минимизируемой нормы). Введём их по~порядку~--- от
одномерного прообраза к~ранг-2 симметричным формам; итоговая сводка
формул приведена в~таблице~\ref{tab:ch1_qn}.

\paragraph{Одномерный прообраз.}
В~скалярном случае ($n=1$) секущий метод заменяет производную разностью
$B_{k+1}=y_k/s_k$, что~однозначно определяет $B_{k+1}$
из~\eqref{eq:secant}. В~$\RR^n$ одного условия недостаточно, и~требуется
дополнительный принцип отбора.

\paragraph{Метод Бройдена.}
Естественный принцип доопределения~--- \emph{минимальное изменение}:
обновление не~должно вносить информацию о~производной в~направлениях,
не~затронутых последним шагом. Формально потребуем, чтобы $B_{k+1}$
действовало как~$B_k$ на~ортогональном дополнении к~$s_k$:
$(B_{k+1}-B_k)\,z = 0$ для~всех $z\perp s_k$. Тогда поправка обращается
в~нуль на~$s_k^\perp$ и~потому имеет ранг~$1$,
$B_{k+1}-B_k = u\,s_k^\top$ с~некоторым $u\in\RR^n$; секущее
условие~\eqref{eq:secant} фиксирует этот вектор:
$u\,\|s_k\|^2 = y_k - B_k s_k$. Обозначив невязку секущей
$\bar y_k = y_k - B_k s_k$, получаем классическое обновление
Бройдена~\cite{broyden1965}
\begin{equation}\label{eq:ch1_broyden}
  B_{k+1} = B_k + \frac{\bar y_k\, s_k^\top}{\|s_k\|^2},
  \qquad \bar y_k = y_k - B_k s_k,
\end{equation}
исторически первое обобщение секущего метода на~многомерный случай.

Тот~же результат отвечает вариационному принципу наименьшего изменения:
\eqref{eq:ch1_broyden}~--- единственное решение задачи
\begin{equation}
  B_{k+1} = \argmin_{B}\ \tfrac12\|B-B_k\|_F^2
  \quad\text{при}\quad B s_k = y_k .
\end{equation}
В~самом деле, лагранжиан
$\mathcal L = \tfrac12\|B-B_k\|_F^2 - \lambda^\top(B s_k - y_k)$
из~условия стационарности даёт $B - B_k = \lambda s_k^\top$,
а~подстановка в~ограничение~--- $\lambda = \bar y_k/\|s_k\|^2$,
что даёт~\eqref{eq:ch1_broyden}; строгая выпуклость нормы
Фробениуса обеспечивает единственность. Геометрически это
ортогональная проекция $B_k$ на~аффинное многообразие
$\{B : B s_k = y_k\}$ в~метрике Фробениуса~\cite[\S8.1]{dennis1996}:
именно в~такой форме обновление переносится на~несколько секущих
(\S\ref{sec:ch1_multisecant}), а~замена нормы Фробениуса на~взвешенную
выделяет симметричные обновления DFP~\eqref{eq:ch1_dfp}
и~BFGS~\eqref{eq:ch1_bfgs}. В данном случае симметрия не~налагается, что~отвечает
общим несимметричным системам~\eqref{eq:system}.

\paragraph{Симметричное ранг-1 обновление (SR1).}
Будем искать симметричную коррекцию ранга~1:
$B_{k+1} = B_k + \delta\,v v^\top$ с~$\delta\in\RR$, $v\in\RR^n$.
Секущее условие~\eqref{eq:secant} даёт
$\delta\,(v^\top s_k)\,v = y_k - B_k s_k = \bar y_k$, откуда $v$
коллинеарен невязке~$\bar y_k$; полагая $v = \bar y_k$, из~равенства
$\delta\,(\bar y_k^\top s_k) = 1$ находим
$\delta = (\bar y_k^\top s_k)^{-1}$ и~обновление
\begin{equation}\label{eq:ch1_sr1}
  B_{k+1} = B_k + \frac{\bar y_k\,\bar y_k^\top}{\bar y_k^\top s_k}.
\end{equation}
Это единственное симметричное ранг-1 обновление, удовлетворяющее
секущей~\cite[\S6.2]{bfgs}. Принцип отбора здесь иной, чем у~метода
Бройдена: симметричным аналогом минимально-фробениусова доопределения
служит не~SR1, а~обновление PSB (Powell symmetric Broyden) ранга~2
(таблица~\ref{tab:ch1_qn})~\cite{powell1970,dennis1977}, тогда как SR1
выделяется требованием минимального ранга коррекции. Тем самым два
естественных принципа отбора в~симметричном классе~--- минимум нормы
Фробениуса и~минимум ранга~--- дают соответственно PSB и~SR1.
Знаменатель $\bar y_k^\top s_k$ может
обращаться в~нуль или быть малым, поэтому~\eqref{eq:ch1_sr1} пропускают
при $|\bar y_k^\top s_k| < r\,\|s_k\|\,\|\bar y_k\|$ ($r\in(0,1)$);
положительная определённость, в~отличие от~BFGS, не~гарантируется, что
компенсируется точностью аппроксимации гессиана в~trust-region
схемах~\cite{conn1991}.

\paragraph{Обновления DFP и BFGS.}
В~оптимизации $B_k$ аппроксимирует гессиан $\nabla^2 f(x_k)$,
и~квазиньютоновский шаг~\eqref{eq:step} есть направление спуска в~точности
тогда, когда $B_k \succ 0$: лишь при этом
$\nabla f(x_k)^\top d_k = -\nabla f(x_k)^\top B_k^{-1}\nabla f(x_k) < 0$.
Это выделяет \emph{симметричные} обновления, к~тому~же сохраняющие
положительную определённость,~--- свойство, которого ни~SR1, ни~метод
Бройдена не~дают. Поэтому к принципу наименьшего изменения добавляются два уточнения:
минимум берётся по~симметричным матрицам, а~норма Фробениуса заменяется на \emph{взвешенную}
\[
  \|A\|_W = \|W^{1/2} A W^{1/2}\|_F, \qquad W \succ 0,
\]
с~метрикой $W$, согласованной с~секущей ($W s_k = y_k$ либо $W y_k = s_k$).
Такую метрику даёт усреднённый гессиан
\[
  \bar G_k = \int_0^1 \nabla^2 f(x_k + \tau s_k)\, d\tau, \qquad
  \bar G_k s_k = y_k.
\]
Существенно, что минимум \emph{не зависит} от~конкретного веса~\cite[\S6.1]{bfgs}, и~потому обновления определены однозначно.

Прямую и~обратную аппроксимации обновляют двойственным образом.
\emph{DFP}~\cite{davidon1991}~--- наименьшее изменение \emph{прямой} матрицы,
\begin{equation}\label{eq:ch1_dfp_var}
  B_{k+1} = \argmin_{B = B^\top}\ \|B - B_k\|_W
  \quad\text{при}\quad B s_k = y_k \quad (W y_k = s_k),
\end{equation}
которое в~привычной записи для~обратной аппроксимации
$H_k = B_k^{-1} \approx [\nabla^2 f(x_k)]^{-1}$ принимает вид
\begin{equation}\label{eq:ch1_dfp}
  H_{k+1} = H_k - \frac{H_k y_k y_k^\top H_k}{y_k^\top H_k y_k}
                 + \frac{s_k s_k^\top}{y_k^\top s_k}.
\end{equation}
Двойственно \emph{BFGS}~--- наименьшее изменение \emph{обратной} матрицы,
\begin{equation}\label{eq:ch1_bfgs_var}
  H_{k+1} = \argmin_{H = H^\top}\ \|H - H_k\|_W
  \quad\text{при}\quad H y_k = s_k \quad (W s_k = y_k),
\end{equation}
что~для~прямой матрицы даёт
\begin{equation}\label{eq:ch1_bfgs}
  B_{k+1} = B_k - \frac{B_k s_k s_k^\top B_k}{s_k^\top B_k s_k}
                 + \frac{y_k y_k^\top}{y_k^\top s_k},
\end{equation}
с~обратным по~формуле Шермана--Моррисона--Вудбури
\begin{equation}\label{eq:ch1_bfgs_inv}
  H_{k+1} = (I - \rho_k s_k y_k^\top)\, H_k\, (I - \rho_k y_k s_k^\top)
            + \rho_k\, s_k s_k^\top, \qquad \rho_k = \frac{1}{y_k^\top s_k}.
\end{equation}
Задачи~\eqref{eq:ch1_dfp_var} и~\eqref{eq:ch1_bfgs_var} переходят друг
в~друга заменой $(B \leftrightarrow H,\ s_k \leftrightarrow y_k)$~--- в~этом
состоит двойственность DFP и~BFGS. Вариационный объект и~привычная форма
записи у~них противоположны: DFP минимально меняет прямую матрицу, но~его
классическая форма~\eqref{eq:ch1_dfp} записана через обратную, а~у~BFGS~---
наоборот.

\paragraph{Положительная определённость.}
Взвешенное наименьшее изменение наследует положительную определённость:
пусть $B_k \succ 0$ и~выполнено условие кривизны
$s_k^\top y_k > 0$~\eqref{eq:curvature}. Тогда для~любого $z \neq 0$
\[
  z^\top B_{k+1} z
   = \Bigl(z^\top B_k z - \frac{(z^\top B_k s_k)^2}{s_k^\top B_k s_k}\Bigr)
     + \frac{(z^\top y_k)^2}{y_k^\top s_k},
\]
где скобка неотрицательна по~неравенству Коши--Буняковского в~скалярном
произведении $\langle a,b\rangle = a^\top B_k b$; второе слагаемое 
$(z^\top y_k)^2/(y_k^\top s_k) > 0$. Т.е. квазиньютоновский шаг остаётся направлением
спуска на~всех итерациях.



\begin{table}[H]
\centering
\caption{Классические квазиньютоновские обновления, удовлетворяющие
секущей $B_{k+1}s_k = y_k$. Через $\bar y_k = y_k - B_k s_k$
обозначена невязка секущей перед обновлением.}\label{tab:ch1_qn}
\renewcommand{\arraystretch}{1.35}
\begin{tabular}{@{}p{3.1cm} p{6.2cm} p{1.6cm} p{1.6cm}@{}}
\toprule
Метод & Формула обновления $B_{k+1} - B_k$ & Симм. & Ранг \\
\midrule
Broyden ("good")~\cite{broyden1965}
  & $\bar y_k\, s_k^\top / \|s_k\|^2$
  & нет & 1 \\
Broyden ("bad")~\cite{broyden1965}
  & $\bar y_k\, y_k^\top / (y_k^\top s_k)$ (для $H_{k+1}$ обратной)
  & нет & 1 \\
SR1~\cite{conn1991,bfgs}
  & $\bar y_k\,\bar y_k^\top / (\bar y_k^\top s_k)$
  & да & 1 \\
PSB~\cite{powell1970,dennis1977}
  & $\dfrac{\bar y_k s_k^\top + s_k \bar y_k^\top}{\|s_k\|^2}
     - \dfrac{(s_k^\top \bar y_k)\, s_k s_k^\top}{\|s_k\|^4}$
  & да & 2 \\
DFP~\cite{davidon1991,dennis1977}
  & ранг-2 формула на $H_{k+1}$, $H_{k+1} y_k = s_k$
  & да & 2 \\
BFGS~\cite{bfgs}
  & $\dfrac{y_k y_k^\top}{y_k^\top s_k} - \dfrac{B_k s_k s_k^\top B_k}{s_k^\top B_k s_k}$
  & да & 2 \\
\bottomrule
\end{tabular}
\end{table}

Параллельно с~плотными обновлениями развивались структурные
модификации: Schubert~\cite{schubert1970} построил версию Бройдена для
систем с разряженным якобианом, сохраняющую заданный шаблон ненулевых
элементов; общая теория ранг-1 обновлений в~семействе
$B_{k+1} = B_k + \bar y_k v_k^\top/(v_k^\top s_k)$ с~произвольным
вектором направления~$v_k$ изложена в~\cite{gay1979}. Эти работы
заложили язык, в~котором обновление параметризуется не~столько
конкретной формулой, сколько выбором матрицы направлений и~базовой
матрицы~--- унифицированную форму этой параметризации
(\S\ref{sec:ch1_multisecant}) систематически развил
Шнабель~\cite{schnabel1983}.


%% ============================================================
\section{Локальная теория сверхлинейной сходимости}\label{sec:ch1_local}
%% ============================================================

Стандартный аппарат локального анализа квазиньютоновских методов
сложился в~работах~\cite{broyden1973,dennis1974,dennis1977,dennis1996}.
Центральный результат Dennis--Mor\'e~\cite{dennis1974}: при
липшицевости якобиана и~невырожденности $J^* = J(x^*)$ итерация
\eqref{eq:step} сходится Q-сверхлинейно тогда~и~только~тогда, когда
выполнено условие
\begin{equation}\label{eq:dm_criterion}
  \lim_{k\to\infty}\frac{\|(B_k - J^*)\,s_k\|}{\|s_k\|} = 0.
\end{equation}
Доказательство сверхлинейности конкретного метода сводится тем~самым
к~проверке~\eqref{eq:dm_criterion}; ключевой инструмент~---
\emph{ограниченная деградация} (bounded
deterioration)~\cite{broyden1973,dennis1996}: при движении вблизи
корня норма ошибки $E_k = B_k - J^*$ может расти не~быстрее, чем
с~линейной добавкой, пропорциональной расстоянию до~корня. Для
классического Бройдена эта оценка получается из~равенства
\begin{equation}\label{eq:bd_rank1}
  \|E_{k+1}\|_F^2 = \|E_k\|_F^2 - \|E_k\,\pi_k\|_F^2 + \Delta_k,
  \qquad \pi_k = \frac{s_k s_k^\top}{\|s_k\|^2},
\end{equation}
где проектор~$\pi_k$~--- ранга~1, а~$\Delta_k = \|y_k-\Jstar s_k\|^2/\|s_k\|^2 = O(\rho_k^2)$
($\rho_k = \max(\|x_k-x^*\|,\|x_{k+1}-x^*\|)$) суммируем при R-линейной
сходимости итераций к~корню. Тождество
\eqref{eq:bd_rank1} вместе с~\eqref{eq:dm_criterion} даёт классическую
теорему Broyden--Dennis--Mor\'e о~локальной Q-сверхлинейности
обновления~\cite{broyden1973}.

Структурная особенность~\eqref{eq:bd_rank1}~--- ранг проектора~$\pi_k$:
отрицательный член разложения <<вычитает>> ошибку только вдоль
направления текущего шага. На~подпространстве $s_k^\perp$ ошибка
переносится из~$E_k$ в~$E_{k+1}$ без~изменений (с~точностью до
липшицевой поправки), что~формально фиксирует
утверждение~\ref{prop:destruction} ниже: ранг-1 обновление разрушает
прошлые секущие на~$s_k^\perp$. Расширение проектора с~ранга~1
за~счёт одновременного сохранения нескольких секущих условий~---
путь усиления этой оценки, разбираемый в~\S\ref{sec:ch1_multisecant}.

Для классических симметричных обновлений локальный анализ устроен
параллельно: Powell~\cite{powell1971} установил Q-сверхлинейность DFP
в~подходящих условиях, конструктивное Dennis--Mor\'e
доказательство~\cite{dennis1977,dennis1996} переносится на~BFGS и~PSB;
для SR1 близкая по~технике локальная теория получена существенно
позже~\cite{conn1991}. Источник дополнительной техники в~симметричном
случае~--- сохранение положительной определённости и~устойчивая работа
оценки $\|E_k\|_F$ под симметричным обновлением.

\paragraph{Современное развитие: явные неасимптотические оценки.}
Классическая теория Денниса--Море характеризует сверхлинейность
\emph{качественно}~--- как асимптотическое свойство~\eqref{eq:dm_criterion},
без~явной скорости. Современная волна работ восполняет этот пробел,
выводя \emph{явные неасимптотические} оценки скорости квазиньютоновских
методов: Rodomanov и~Nesterov установили явную сверхлинейную сходимость
жадных (greedy) обновлений~\cite{rodomanov2021greedy} и~неасимптотические
оценки для~классических обновлений~\cite{rodomanov2022rates}; Jin
и~Mokhtari~\cite{jin2023nonasymp} и~Lin, Ye, Zhang~\cite{lin2022explicit}
получили явные оценки скорости для~стандартных, жадных
и~рандомизированных схем. Эти результаты делают для~\emph{скорости
сходимости} то~же, что настоящая работа~--- для~\emph{оценки ограниченной
деградации}: переводят качественную картину в~количественную с~явными
константами. Развиваемое в~гл.~\ref{ch:sp-broyden-theory} точное
разложение нормы ошибки с~явной геометрической константой
$C_{\mathrm{PB}}$ примыкает к~этой программе, отличаясь объектом:
мультисекущее обновление для~нелинейных систем и~количественная оценка
самой ошибки $\|B_k-J^*\|_F$, а~не~асимптотической скорости.

%% ============================================================
\section{Мультисекущие обновления: общая формула и таксономия}\label{sec:ch1_multisecant}
%% ============================================================

Пусть заданы $m$ секущих условий $B_{k+1} S_m = Y_m$, где
\[
  S_m = [s_{k-m+1}\mid\cdots\mid s_k] \in \RR^{n \times m}, \qquad
  Y_m = [y_{k-m+1}\mid\cdots\mid y_k] \in \RR^{n \times m}.
\]


Семейство обновлений, удовлетворяющих всем $m$ секущим условиям,
описывается унифицированной формулой
Шнабеля~\cite[ур.~(2.5)]{schnabel1983}:
\begin{equation}\label{eq:general}
  \boxed{
    B_{k+1} = \widehat{B}_k + (Y_m - \widehat{B}_k S_m)\,(V_m^\top S_m)^{-1}\,V_m^\top
  }
\end{equation}
где $\widehat{B}_k \in \RR^{n \times n}$~--- базовая матрица,
$V_m \in \RR^{n \times m}$~--- матрица направлений
($V_m^\top S_m$ невырождена). Прямая подстановка даёт
$B_{k+1} S_m = Y_m$ при~любом допустимом~$V_m$; ранг
$\Delta_k = B_{k+1}-\widehat{B}_k$ не~превосходит~$m$.

\subsection*{Частный случай: $m = 1$}

При $m = 1$ имеем $S_m = s_k$, $Y_m = y_k$, $V_m = v_k \in \RR^n$,
и формула~\eqref{eq:general} принимает вид
\begin{equation}\label{eq:rank1}
  B_{k+1} = \widehat{B}_k + \frac{(y_k - \widehat{B}_k s_k)\,v_k^\top}{v_k^\top s_k}.
\end{equation}
Это ранг-1 обновление, удовлетворяющее единственному секущему условию
$B_{k+1} s_k = y_k$. При дополнительном ограничении $v_k^\top s_k = 1$
знаменатель исчезает:
\begin{equation}\label{eq:rank1_normalized}
  B_{k+1} = \widehat{B}_k + (y_k - \widehat{B}_k s_k)\,v_k^\top.
\end{equation}

\subsection*{Разрушение прошлых секущих ранг-1 обновлением}\label{sec:rank1_destruction}

Ранг-1 ветвь~\eqref{eq:rank1} обладает характерной патологией:
текущее секущее условие $B_{k+1}s_k = y_k$ выполняется по~построению, однако
\emph{прошлые} секущие в~общем случае нарушаются. Это и~есть основная
мотивация перехода от~классической ранг-1 формулы к~мультисекущим
расширениям~\cite{schnabel1983,gay_schnabel1978}, рассматриваемым
в~гл.~\ref{ch:sp-broyden-theory}.

\begin{proposition}\label{prop:destruction}
Пусть $B_{k+1}$ задано ранг-1 обновлением~\eqref{eq:rank1} в~режиме
накопления $\widehat{B}_k = B_k$ с~произвольным $v_k$ ($v_k^\top s_k \neq 0$),
\[
  B_{k+1} = B_k + \frac{(y_k - B_k s_k)\,v_k^\top}{v_k^\top s_k}.
\]
Тогда для любого $j < k$:
\[
  B_{k+1} s_j = B_k s_j + (y_k - B_k s_k)\,\frac{v_k^\top s_j}{v_k^\top s_k}.
\]
В частности, если $y_k \neq B_k s_k$ (т.\,е.\ текущая секущая ещё не~выполнена),
то $B_{k+1} s_j = B_k s_j$ тогда и~только тогда, когда $v_k^\top s_j = 0$.
\end{proposition}

\begin{proof}
Прямое вычисление даёт
\[
  B_{k+1} s_j
  = B_k s_j + (y_k - B_k s_k)\,\frac{v_k^\top s_j}{v_k^\top s_k}.
\]
При $y_k \neq B_k s_k$ второе слагаемое зануляется тогда~и~только тогда,
когда $v_k^\top s_j = 0$.
\end{proof}

\begin{corollary}\label{cor:broyden_destruction}
В~классическом методе Бройдена ($v_k = s_k$), если на~предыдущем шаге
выполнялось $B_k s_j = y_j$ и~$y_k \neq B_k s_k$, то $B_{k+1} s_j = y_j$
тогда~и~только тогда, когда $s_k^\top s_j = 0$. В~общем случае равенство нарушается.
\end{corollary}

Естественный способ устранить это нарушение~--- потребовать сохранения
\emph{всех} прошлых секущих в~окне: $B_{k+1} s_j = y_j$ для
$j=k{-}p,\ldots,k$. Это и~приводит к~мультисекущему обновлению
Бройдена~\cite{schnabel1983} (гл.~\ref{ch:sp-broyden-theory}) и~к~его
симметричному аналогу~SS-$\mathcal{U}$ (гл.~\ref{ch:symmetric}); именно
сохранение прошлых секущих расширяет проектор bounded deterioration
с~ранга~1 до~ранга~$p{+}1$, что~и~составляет центральный объект
анализа последующих глав.

\subsection*{Выбор $V_m$: что минимизирует обновление}

При фиксированных $\widehat{B}_k$, $S_m$, $Y_m$ выбор $V_m$ определяет
\emph{форму} обновления $\Delta_k$. Различные $V_m$ удовлетворяют
одним и тем же секущим условиям, но отличаются по норме изменения
$\|\Delta_k\|_F = \|B_{k+1} - \widehat{B}_k\|_F$.

\paragraph{$V_m = S_m$: минимум нормы Фробениуса.}
Обновление принимает вид
\begin{equation}\label{eq:Vm_Sm}
  B_{k+1} = \widehat{B}_k + (Y_m - \widehat{B}_k S_m)\,(S_m^\top S_m)^{-1}\,S_m^\top.
\end{equation}
Это единственное решение задачи $\min_{B}\|B - \widehat{B}_k\|_F$
при $B\,S_m = Y_m$~--- проекция $\widehat{B}_k$ на~аффинное подпространство,
задаваемое секущими условиями, в~норме Фробениуса; строгая выпуклость
гарантирует единственность. Ранг обновления не~превосходит~$m$.

\paragraph{$m=1$, $v_k = s_k$: классический метод Бройдена.}
Частный случай при $m=1$:
\begin{equation}\label{eq:good_broyden}
  B_{k+1} = \widehat{B}_k + \frac{(y_k - \widehat{B}_k s_k)\,s_k^\top}{\|s_k\|^2}.
\end{equation}
Это решение $\min_B \|B - \widehat{B}_k\|_F$ при $B\,s_k = y_k$;
при $\widehat{B}_k = B_k$~--- классический метод
Бройдена~\cite{broyden1965}, ранг~1, $O(n^2)$ на шаг. Другие
варианты~--- ``bad Broyden'' ($v_k=y_k$), отвечающий
$\min\|H_{k+1}-H_k\|_F$ при $H_{k+1} y_k=s_k$ для обратного
якобиана~\cite{broyden1965,dennis1996}, и~произвольный $v_k$
с~$v_k^\top s_k\ne 0$~\cite{gay1979}~--- в~настоящей работе
не~используются и~упоминаются для полноты семейства~\eqref{eq:rank1}.


%% ============================================================
\section{Базовая матрица и связь с ускорением Андерсона}\label{sec:ch1_base}
%% ============================================================

Второй параметр семейства~\eqref{eq:general}~--- базовая
матрица~$\widehat{B}_k$. При $\widehat{B}_k=B_k$ (\emph{накопление})
вся прошлая информация сохраняется неявно в~$B_k$; этот режим~---
стандартный для метода Бройдена и~основной в~настоящей работе.


При $\widehat{B}_k=B_0$ (\emph{пересборка}, обычно $B_0=I$) обновление
зависит только от~$B_0$ и~окна~$m$ пар, история за~его пределами
стирается. В~частности, при $V_m=S_m$ и~$B_0=I$
формула~\eqref{eq:Vm_Sm} принимает вид
\begin{equation}\label{eq:anderson}
  B_{k+1} = I + (Y_m - S_m)\,(S_m^\top S_m)^{-1}\,S_m^\top,
\end{equation}
и~шаг $x_{k+1}=x_k - B_{k+1}^{-1} F(x_k)$ совпадает с~ускорением
Андерсона типа~I при коэффициенте релаксации
$\beta=1$~\cite{anderson1965,walker2011,fang2009}. Современную
таксономию мультисекущих ускорителей развивают
Fang--Saad~\cite{fang2009}, выделяя два класса (Type-I/Type-II),
эквивалентных, соответственно, Бройдену и~Андерсону, а~локальную
теорию сходимости Андерсона~--- Toth--Kelley~\cite{toth2015}; глобально
сходящийся вариант type-I (бройденовского типа) для~негладких задач
о~неподвижной точке развит Zhang, O'Donoghue и~Boyd~\cite{zhang2020anderson}. В~режиме
накопления, основном для гл.~\ref{ch:sp-broyden-theory}, обновление
наследует поведение Бройдена с~расширенным окном; в~Алгоритме
гл.~\ref{ch:sp-broyden-theory} режим пересборки активируется опцией
$\mathtt{accumulate}=\text{false}$ и~используется в~численных
экспериментах для сравнения.


%% ============================================================
\section{Глобализация и limited-memory варианты}\label{sec:ch1_global}
%% ============================================================

Локальная теория \S\ref{sec:ch1_local} гарантирует сверхлинейную
сходимость из~достаточно малой окрестности корня~$x^*$. Для практического применения нужна
\emph{глобализация}: стратегия, обеспечивающая сходимость к~корню
из~произвольной начальной точки в~разумной области сходимости. Параллельная
задача в~крупномасштабном режиме~--- \emph{limited-memory} реализация,
позволяющая работать с~$B_k$, не~храня плотной матрицы $n\times n$.

\subsection{Глобализация для нелинейных систем и оптимизации}

В~гладкой оптимизации глобализация классических квазиньютоновских
методов опирается на~монотонное убывание целевой функции $f$ и~линейный
поиск, удовлетворяющий условиям Армихо~\cite{armijo1966} или
Вольфа~\cite{wolfe1969}; общая теория глобальной сходимости
направления спуска восходит к~Zoutendijk~\cite{zoutendijk1970}.
Сохранение положительной определённости BFGS под условием Вольфа даёт
сходимость на~достаточно широком классе задач, тогда~как
для PSB и~SR1 теория глобальной сходимости заметно слабее: их обновления
не~сохраняют знакоопределённости, что~требует trust-region схем
или гибридизации. Альтернативный подход к~глобализации для
оптимизации~--- регуляризация (Levenberg--Marquardt~\cite{marquardt1963})
в~приложении к~нелинейным наименьшим квадратам; в~общем оптимизационном
контексте регуляризация естественно проявляется как trust-region
ограничение на~шаг.

Для нелинейных систем $F(x)=0$ ситуация существенно отличается:
монотонной функции, минимум которой автоматически соответствует
корню, нет. Классический эмпирический
анализ показывает, что~даже на~хорошо обусловленных задачах
поведение метода Бройдена чувствительно к~начальному
приближению~\cite{shannophua1978,more1981}. Стандартные стратегии
глобализации в~этом классе~--- линейный поиск по~$\|F\|$ с~условием
типа Армихо и~неточные ньютоновские шаги (inexact Newton).
Нестандартный, но~важный для настоящей работы инструмент~---
\emph{немонотонный} линейный поиск, допускающий временный рост
невязки~$\|F\|$ при условии возврата к~монотонной траектории через
скользящее среднее или максимум по~окну предыдущих значений; такой
поиск, в~частности, использован Бурдаковым и~Каманди~\cite{burdakov2018}
в~глобализованном multipoint secant методе для $F(x)=0$. Их работа~---
ближайший прецедент глобализации мультисекущих обновлений для
нелинейных систем; ключевое отличие от~режима, рассматриваемого
в~настоящей работе,~--- в~технике переноса bounded deterioration
через вспомогательную аппроксимацию якобиана на~шаге $k+1$,
что~не~требует повторного анализа разложения~\eqref{eq:bd_rank1}
для каждого варианта линейного поиска.

\subsection{Limited-memory реализации}

Для крупномасштабных задач хранение плотного $B_k\in\RR^{n\times n}$
($O(n^2)$ памяти) затруднительно уже при $n\sim 10^5$. Канонический
limited-memory вариант квазиньютоновского метода для оптимизации~---
L-BFGS ~\cite{lbfgs}: вместо плотной $B_k$ хранятся
последние $m$ пар $(s_j, y_j)$, через которые $B_k v$ или $H_k v$
восстанавливаются процедурой two-loop рекурсии за~$O(mn)$ операций.
Liu и~Nocedal~\cite{liu1989} продемонстрировали практическую
эффективность L-BFGS на~крупномасштабных задачах; Byrd, Nocedal
и~Schnabel~\cite{byrd1994} получили компактное матричное
представление L-BFGS и~L-SR1, удобное для trust-region схем
и~теоретического анализа.

Для несимметричных задач (нелинейные системы~\eqref{eq:system})
канонической limited-memory формы классического Бройдена в~литературе
\emph{не~сложилось}: прямое сокращение хранения через окно $m$ пар
$(s_j, y_j)$ ведёт к~потере секущих условий вне окна (ср.\
утв.~\ref{prop:destruction}), и~теория ограниченной деградации
для такой <<скользящей>> версии Бройдена существенно слабее, чем
для L-BFGS. Мультисекущая формулировка~\eqref{eq:general} в~режиме
накопления естественно даёт ранг-$(p{+}1)$ обновление, согласованное
сразу со~всеми парами окна, что~открывает путь к~limited-memory
варианту с~теоретическими гарантиями, не~деградирующими с~ростом~$n$.
Этот путь реализуется в~гл.~\ref{ch:sp-broyden-theory},
\S\ref{sec:highdim} в~виде L-PB.


%% ============================================================
\section{Симметричные multi-secant подходы}\label{sec:ch1_sym}
%% ============================================================

Перенос идеи сохранения нескольких секущих условий на~симметричный
случай (для гладкой оптимизации) наталкивается на~классическое
препятствие, явно зафиксированное у~Шнабеля~\cite{schnabel1983}:
система $B_{k+1} S_m = Y_m$ при $B_{k+1} = B_{k+1}^\top$ совместна
тогда и~только тогда, когда $S_m^\top Y_m$ симметрична, что~в~общем
случае не~выполняется. Дословный перенос мультисекущей формулы
с~симметризацией невозможен, и~литература последних лет предлагает
различные стратегии \emph{компромисса} между симметрией и~точным
выполнением прошлых секущих условий.

\begin{itemize}[leftmargin=*]
  \item Gratton, Malmedy, Toint~\cite{gratton2015}: \emph{weighted
        secant equations}~--- секущие условия по~прошлым парам
        выполняются приближённо с~явными весами, симметрия
        обновления сохраняется; задача сводится к~взвешенной
        фробениусовой проекции с~квадратичным критерием.
  \item Boutet, Haelterman, Degroote~\cite{boutet2020,boutet2021,boutet2022}:
        семейство симметричных PSB-расширений (SUpPSB, gPSB,
        grouped/ordered PSB), различающихся способом
        регуляризации совместности; общая черта~--- замена
        точного выполнения прошлых секущих на~некоторое
        статистическое или групповое условие, сохраняющее
        симметрию ранг-2 коррекции.
  \item Scieur, Liu, Pumir, Boumal~\cite{scieur2021}: regularized
        least-squares форма с~явной регуляризацией грамиана и~общим
        обобщённым обоснованием как для несимметричного, так и~для
        симметричного случая; точное выполнение секущих условий
        приносится в~жертву устойчивости при плохо обусловленном
        окне.
\end{itemize}

Каждый из~перечисленных подходов сохраняет симметрию (а~часто
и~положительную определённость) ценой ослабления секущих условий~---
текущего, прошлых или того и~другого. В~гл.~\ref{ch:symmetric}
рассматривается альтернативная стратегия компромисса
SS-$\mathcal{U}$: точное выполнение \emph{текущей} секущей
$B_{k+1}s_k = y_k$ и~симметрии, а~прошлые секущие стабилизируются
ранг-1 коррекцией в~подпространстве $s_k^\perp$; детальное сравнение
SS-$\mathcal{U}$ с~работами~\cite{gratton2015,boutet2020,boutet2021,boutet2022,scieur2021}
проведено в~\S\ref{sec:ss_related}.


%% ============================================================
\section{Место настоящей работы}\label{sec:ch1_place}
%% ============================================================

Объект анализа настоящей диссертации~--- мультисекущее обновление
Бройдена в~форме~\eqref{eq:general} при $\widehat{B}_k = B_k$, $V_m = S_m$,
$m = p+1$ (режим накопления, минимум-Фробениус); этот объект введён
Шнабелем~\cite{schnabel1983}, а~ранг-1 специализация под индуктивным
инвариантом $B_k s_j = y_j$ для~$j<k$~--- Гэем и
Шнабелем~\cite{gay_schnabel1978}. Развиваемое в~гл.~\ref{ch:sp-broyden-theory}
направление: количественная форма теоремы об ограниченной деградации
с~явной константой через геометрию окна $\kappa_p = \kappa(S_p)$,
точное разложение нормы ошибки вида~\eqref{eq:bd_rank1}, но~с~проектором
ранга~$p{+}1$ (вместо ранга~1 у~классического Бройдена), $O(\rho_k)$-аппроксимация
якобиана на~подпространстве прошлых шагов, а~также limited-memory
реализация L-PB через ранг-1 запись Гэя--Шнабеля.

Для симметричного случая (гл.~\ref{ch:symmetric}) развивается
секантно-стабилизированное обновление~SS-$\mathcal{U}$~--- симметричный
аналог сохранения прошлых секущих, разрешающий классическую
несовместность симметрии и~сохранения нескольких секущих
условий~\cite{schnabel1983} ранг-1 коррекцией в~$s_k^\perp$ (механизм
и~сравнение с~известными подходами~--- \S\ref{sec:ch1_sym}).
Установлена Q-сверхлинейная сходимость~SS-$\mathcal{U}$ по~критерию
Денниса--Море.

\begin{table}[ht]
\centering
\caption{Сводка оценок ограниченной деградации и~место настоящей работы:
классический и~projected Бройден, мультисекущее обновление
(гл.~\ref{ch:sp-broyden-theory}), блочный Бройден и~современная волна
явных оценок скорости.}\label{tab:ch1_overview}
\renewcommand{\arraystretch}{1.35}
\begin{tabular}{|@{}p{2.9cm}|p{4.0cm}|p{3.8cm}|p{4.3cm}@{}|}
\toprule
Метод & Проектор / ранг & Характер оценки $\|E_{k+1}\|_F^2$ & Константа / класс \\
\midrule
Классический Бройден~\cite{broyden1965,dennis1996}
  & $s_k s_k^\top/\|s_k\|^2$, ранг~1
  & равенство (Lemma~8.2.1), контроль вдоль~$s_k$
  & $L(1{+}C_s/2)$;\newline $F(x)=0$ \\ \hline
Projected Broyden~\cite{gay_schnabel1978}
  & ранг~1 (инвариант $B_k s_j{=}y_j$)
  & оценка \emph{сверху} (Th~4.2)
  & ---;\newline $F(x)=0$ \\ \hline
Мультисекущий \emph{(гл.~\ref{ch:sp-broyden-theory})}
  & $S_p G_p^{-1} S_p^\top$, ранг~$p{+}1$
  & точное \emph{равенство} \eqref{eq:exact_decomp}, $O(\rho_k)$ на~$\mathrm{col}\,S_p$
  & $C_{\mathrm{PB}}{=}L(1{+}\tfrac{C_s}{2})\sqrt{p{+}1}\,\kappa_p$;\newline $F(x)=0$ \\ \hline
Block Broyden~\cite{liu2023block}
  & $U(U^\top U)^{-1}U^\top$ случ., ранг~$k$
  & равенство \emph{в~среднем} $(1{-}k/d)$
  & condition-number-free;\newline $F(x)=0$ \\ \hline
Явные оценки скорости~\cite{rodomanov2021greedy,rodomanov2022rates,jin2023nonasymp,lin2022explicit}
  & --- (greedy / random)
  & оценки \emph{скорости}, не~BD
  & $\dim\times\log\kappa$;\newline $\min f$ \\
\bottomrule
\end{tabular}
\end{table}

\paragraph{Резюме.}
Общая мультисекантная формула~\eqref{eq:general} параметризуется
матрицей направлений~$V_m$ (форма обновления: минимум по~Фробениусу
при~$V_m = S_m$, классический Бройден при $m=1$, $v_k = s_k$;
\eqref{eq:Vm_Sm}, \eqref{eq:good_broyden}) и~базовой
матрицей~$\widehat{B}_k$ (накопление или пересборка, последняя
эквивалентна Андерсону;~\eqref{eq:anderson}). Локальная теория
сверхлинейной сходимости~\cite{broyden1973,dennis1974,dennis1977}
сводит вопрос к~критерию Денниса--Море~\eqref{eq:dm_criterion}
и~разложению bounded deterioration~\eqref{eq:bd_rank1} с~проектором
ранга~1 для классических обновлений; разрушение прошлых секущих
ранг-1 обновлением (утв.~\ref{prop:destruction}) мотивирует переход
к~мультисекущим расширениям. Дальнейший анализ ведётся в~режиме
накопления с~расширенным окном секущих~--- мультисекущим обновлением
Бройдена~\cite{schnabel1983} (гл.~\ref{ch:sp-broyden-theory})
и~его симметричным аналогом
SS-$\mathcal{U}$ (гл.~\ref{ch:symmetric}).


%% ============================================================
